\documentclass[10pt,a4paper]{amsart}
\usepackage[margin=19mm]{geometry}
\usepackage{amsmath,amssymb,mathtools}
\usepackage[scr=rsfs,cal=euler]{mathalfa}
\usepackage{xfrac}
\usepackage[unicode,hidelinks,bookmarks=false]{hyperref}
\newcommand{\ie}{i.e.\ }
\newcommand{\cf}{cf.\ }
\newcommand{\resp}{respectively\ }
\newcommand{\Subsection}{\subsection}
\providecommand{\nobreakdash}{\nobreak\hbox{-}}
\setlength{\emergencystretch}{2em}
\multlinegap0pt
\usepackage{amssymb}
\usepackage{dsfont}
\usepackage{xcolor}
\usepackage{cancel}
\usepackage[shortlabels]{enumitem}
\usepackage[normalem]{ulem}
\newtheorem{proposition}{}
\newcommand\Per{\mathop{\rm Per}}
\def\XXint#1#2#3{{\setbox0=\hbox{$#1{#2#3}{\int}$}
  \vcenter{\hbox{$#2#3$}}\kern-.5\wd0}}
\def\Xint#1{\mathchoice{\XXint\displaystyle\textstyle{#1}}%
  {\XXint\textstyle\scriptstyle{#1}}%
  {\XXint\scriptstyle\scriptscriptstyle{#1}}%
  {\XXint\scriptscriptstyle\scriptscriptstyle{#1}}%
  \!\int}
\newcommand{\dHn}{\,\dd\mathcal{H}^{n-1}}
\def\dd{{\rm d}}
\def\ddiv{{\mathop{\rm div}}}
\newcommand{\eqdef}{\overset{\text{def.}}{=}}
\def\fint{\Xint-}
\title{On the effective sharp stability of the Faber-Krahn inequality}
\author{Andr\'e Guerra, Jo\~ao Miguel Machado, Jo\~ao P. G. Ramos}
\date{Source: arXiv:2609.07772v1 (CC BY 4.0)}
\begin{document}
\maketitle

\noindent\emph{Source and licence.} On the effective sharp stability of the Faber-Krahn inequality. Version arXiv:2609.07772v1 (CC BY 4.0), distributed by arXiv under the Creative Commons Attribution 4.0 International licence, http://creativecommons.org/licenses/by/4.0/, as recorded in the arXiv licence metadata for this exact version and shown on the abstract page https://arxiv.org/abs/2609.07772v1. Full licence notice: \texttt{licenses/arxiv-2609.07772-CC-BY-4.0.txt}.\par\smallskip

\noindent\textit{Editorial scope.} Counted unit: the article's proposition prop.estimates\_q'\_DH\_DI (source0.tex lines 1606--1640). The article's complete original proof of it is delivered with it (source0.tex lines 1642--1885, one \texttt{proof} environment of 244 lines). No mathematics is written, repaired or re-derived: the statement and the proof are the article's own lines, in the article's own notation, and no retained line is substituted or re-flowed.  Retained uncounted as the source-local context the proof needs: lines 688--694; lines 922--926; lines 1523--1539.  Nothing was omitted from any retained span, so the package carries no omission record.  No retained line contains a citation, so the wrapper carries no bibliography and no bibliography entry.  No retained line contains a figure environment, an image-inclusion command or drawing code, so no image is delivered and the package declares no asset dependency.  Editorial formatting only, in addition: an AMS \texttt{amsart} preamble that restates the article's own declarations and macros used by the retained lines, namely the environments \texttt{proposition} on separate counters, together with the standard packages those lines need.  The retained environments print in this document as Proposition 1 in order of appearance, whereas the article prints the same items with the numbering of its own full document.  The complete native source of the article as distributed in the arXiv e-print for this exact version is bundled unchanged as \texttt{sources/209662/source0.tex}.
\begin{equation}\label{eq.quant_isoperimetric}
  \Per(E)-\Per(B_{r_E})
  \ge \kappa_n r_E^{n-1}\mathcal A(E)^2,
  \qquad
  \mathcal A(E)\eqdef
  \inf_{z\in\mathbb R^n}\frac{|E\Delta B_{r_E}(z)|}{|E|}.
\end{equation}

\begin{equation}\label{eq.velocity_normals}
    V_\rho\eqdef\frac{w(\rho)}{|\nabla u|},\qquad
    H_{z,\rho}(x)\eqdef\frac{(x-z)\cdot\nu_\rho(x)}{\rho},\qquad
    \overline{V}_\rho\eqdef\fint_{\partial^*E_\rho}V_\rho\ =\ \frac{\Per(B_\rho)}{\Per(E_\rho)},
\end{equation}

 \begin{proposition}
 \label{prop.FJ-strong}
For every   $n\ge2$ and $R_0\ge2$ there are computable constants
$\varepsilon_{\rm FJ}(n,R_0)>0$ and $C_{\rm FJ}(n,R_0)>0$ such that every  set
of finite perimeter   $F\subset B_{R_0}$ satisfying
\[
 |F|=\omega_n,\qquad \int_Fx\dd x=0,
 \qquad \Per(F)-\Per(B_1)\le\varepsilon_{\rm FJ}(n,R_0)
\]
obeys
\begin{equation}\label{eq.FJ-exact}
 |F\Delta B_1|^2
 +\int_{\partial^*F}
 \left|\nu_F-\frac{x}{|x|}\right|^2\dHn
 \le C_{\rm FJ}(n,R_0)\bigl(\Per(F)-\Per(B_1)\bigr).
\end{equation}
\end{proposition}

\begin{proposition}\label{prop.estimates_q'_DH_DI}
The following estimates hold for almost every $\rho\in(0,1)$:
\begin{itemize}
\item[(1)]
  \[
     \left(\int_{\partial^*E_\rho}
    |V_\rho-\overline V_\rho|\dHn\right)^2
    \le D_H(t(\rho)).
 \]
 \item[(2)] With the computable constant $\kappa_n$ from
\eqref{eq.quant_isoperimetric},
\[
 q(\rho)^2
 \le\frac{\omega_n}{2n\kappa_n}\rho^2D_I(t(\rho)).
\]
\item[(3)] For every $\rho_*\in(0,1)$ there are computable constants
$\eta_I=\eta_I(n,R,\rho_*)>0$ and $C=C(n,R,\rho_*)>0$ such that, for almost
every $\rho\in[\rho_*,1]$ satisfying
$D_I(t(\rho))\le\eta_I$, there is a center $\bar z_\rho$ for which
\[
 |E_\rho\Delta B_\rho(\bar z_\rho)|^2
 +\left(\int_{\partial^*E_\rho}
 |H_{\bar z_\rho,\rho}-\overline V_\rho|\dHn\right)^2
 \le C D_I(t(\rho)).
\]
If  $z_\rho$ is optimal for $q(\rho)$, then also
\[
 |z_\rho-\bar z_\rho|\le C D_I(t(\rho))^{1/2},
 \qquad
 \left(\int_{\partial^*E_\rho}
 |H_{z_\rho,\rho}-\overline V_\rho|\dHn\right)^2
 \le C D_I(t(\rho)).
\]
 \end{itemize}
\end{proposition}

\begin{proof}
We prove the three estimates at radii for which the preceding identities hold and $E_\rho$ has finite perimeter. Starting with item (1), recall the definition of $V_\rho$ and $\overline{V}_\rho$ in~\eqref{eq.velocity_normals}
\[
  V_\rho\eqdef\frac{w(\rho)}{|\nabla u|},\qquad
  \overline{V}_\rho\eqdef\fint_{\partial^*E_\rho}V_\rho\dHn=\frac{\Per(B_\rho)}{\Per(E_\rho)},
\]
as well as the deficit $D_H(t(\rho))=\alpha(t(\rho))\,\beta(t(\rho))-\Per(E_\rho)^2$. By the Cauchy-Schwarz inequality, it holds that
\begin{align*}
  \left(\int_{\partial^*E_\rho}
  \bigl|V_\rho-\overline{V}_\rho\bigr|\dHn\right)^{2}
  \le
  \int_{\partial^*E_\rho}V_\rho\dHn
  \int_{\partial^*E_\rho}\frac{\bigl(V_\rho-\overline{V}_\rho\bigr)^2}{V_\rho}\dHn
\end{align*}
The first term can be computed as
\[
  \int_{\partial^*E_\rho}V_\rho\dHn = \Per(E_\rho) \bar V_\rho = \Per(B_\rho).
\]
As for the second, we develop the squares to obtain
  \begin{align*}
  \int_{\partial^*E_\rho}\frac{\bigl(V_\rho-\overline{V}_\rho\bigr)^2}{V_\rho}\dHn
  = &
  w(\rho)\int_{\partial^* E_\rho} \frac{1}{|\nabla u|} \dHn\\
  & 
  -2\bar V_\rho \Per(E_\rho)
  +
  \frac{\bar V_\rho^2}{w(\rho)} \int_{\partial^* E_\rho}|\nabla u| \dHn
\end{align*}
Hence, recalling the definitions of $\alpha(t), \beta(t)$ and that $w(\rho)\alpha(t(\rho))=\Per(B_\rho)$ we get
\begin{align*}
    \int_{\partial^*E_\rho}\frac{\bigl(V_\rho-\overline{V}_\rho\bigr)^2}{V_\rho}\dHn
    &=\frac{\Per(B_\rho)}{\Per(E_\rho)^2} \alpha(t(\rho))\,\beta(t(\rho))-\Per(B_\rho)
    =\frac{\Per(B_\rho)}{\Per(E_\rho)^2} D_H(t(\rho)).
\end{align*}
Multiplying both terms, and recalling that $\Per(B_\rho) \le \Per(E_\rho)$
by the isoperimetric inequality, since these sets have the same volume by construction, item (1) follows.

For item~(2), the quantitative isoperimetric inequality gives
\begin{align*}
 D_I(t(\rho))
 &\ge 2n\omega_n\kappa_n\rho^{2n-2}\mathcal A(E_\rho)^2
 =\frac{2n\kappa_n}{\omega_n\rho^2}q(\rho)^2,
\end{align*}
which is the asserted estimate.

We now prove item~(3). Fix $\rho_*\in(0,1)$ and
$\rho\in[\rho_*,1]$. Let
\[
 \bar z_\rho\eqdef\frac{1}{|E_\rho|}\int_{E_\rho}x\dd x
\]
be the barycenter of $E_\rho$, and define
$F_\rho=\rho^{-1}(E_\rho-\bar z_\rho)$. Since
$E_\rho\subset B_R$, one has $|\bar z_\rho|\le R$ and
$F_\rho\subset B_{2R/\rho_*}$. Moreover,
\begin{align}\label{eq.normalized_perimeter_deficit}
 \Per(F_\rho)-\Per(B_1)
 &=\frac{\Per(E_\rho)-\Per(B_\rho)}{\rho^{n-1}}
  \le\frac{D_I(t(\rho))}{2n\omega_n\rho^{2n-2}}
 \le\frac{D_I(t(\rho))}{2n\omega_n\rho_*^{2n-2}}.
\end{align}
Choose the computable number $\eta_I(n,R,\rho_*)\le1$ so that the last
quantity is at most
$\varepsilon_{\rm FJ}(n,\max\{2,2R/\rho_*\})$ whenever
$D_I(t(\rho))\le\eta_I$. Proposition~\ref{prop.FJ-strong}, followed by
scaling, then gives
\begin{align}
 |E_\rho\Delta B_\rho(\bar z_\rho)|^2
 &\le C(n,R,\rho_*)D_I(t(\rho)),
 \label{eq.local_FJ_volume}\\
 \int_{\partial^*E_\rho}
 \left|\nu_\rho-\frac{x-\bar z_\rho}{|x-\bar z_\rho|}\right|^2\dHn
 &\le C(n,R,\rho_*)D_I(t(\rho)).
 \label{eq.local_FJ_normal}
\end{align}

We first estimate the second quantity in item~(3). Recall that
$
 \overline V_\rho=\frac{\Per(B_\rho)}{\Per(E_\rho)}.
$
The triangle inequality gives
\begin{align*}
 \int_{\partial^*E_\rho}
 |H_{\bar z_\rho,\rho}-\overline V_\rho|\dHn
 &\le \Per(E_\rho)(1-\overline V_\rho)
 +\int_{\partial^*E_\rho}|H_{\bar z_\rho,\rho}-1|\dHn.
\end{align*}
The first term above, which takes into account the error of replacing
$\overline V_\rho$ by $1$, satisfies
\begin{align*}
 \Per(E_\rho)(1-\overline V_\rho)
 &=\Per(E_\rho)-\Per(B_\rho)
 \le\frac{D_I(t(\rho))}{2n\omega_n\rho^{n-1}}
 \le C(n,\rho_*)D_I(t(\rho)).
\end{align*}

As for the second term, set
$
 a\eqdef\frac{|x-\bar z_\rho|}{\rho}
$,
$
 b\eqdef\frac{x-\bar z_\rho}{|x-\bar z_\rho|}\cdot\nu_\rho,
$
so that $H_{\bar z_\rho,\rho}=ab$, $|b|\le1$, and
\begin{align*}
 |H_{\bar z_\rho,\rho}-1|
 &=|ab-1|=|(a-1)b+(b-1)|
 \le |a-1||b|+1-b\\
 &\le\left|\frac{|x-\bar z_\rho|}{\rho}-1\right|
 +\frac12\left|\nu_\rho-
 \frac{x-\bar z_\rho}{|x-\bar z_\rho|}\right|^2.
\end{align*}
Here we used that both vectors in the last line have unit length, and hence
$1-b$ is one half of their squared distance. Consequently,
\[
 \int_{\partial^*E_\rho}|H_{\bar z_\rho,\rho}-1|\dHn
 \le \mathcal I_1+\mathcal I_2,
\]
where
\begin{align*}
 \mathcal I_1\eqdef\int_{\partial^*E_\rho}|a-1|\dHn, \quad 
 \mathcal I_2\eqdef\frac12\int_{\partial^*E_\rho}
 \left|\nu_\rho-\frac{x-\bar z_\rho}{|x-\bar z_\rho|}\right|^2\dHn.
\end{align*}
Estimate~\eqref{eq.local_FJ_normal} directly gives
\begin{equation}\label{eq.I2_estimate}
 \mathcal I_2\le C(n,R,\rho_*)D_I(t(\rho)).
\end{equation}

For $\mathcal I_1$, we use again the identity
$\displaystyle
 1=b+\frac12\left|\nu_\rho-
 \frac{x-\bar z_\rho}{|x-\bar z_\rho|}\right|^2
$
to write
\begin{align*}
 \mathcal I_1
 &=\underbrace{\int_{\partial^*E_\rho}|a-1|
 \frac{x-\bar z_\rho}{|x-\bar z_\rho|}\cdot\nu_\rho\dHn}
 _{\eqdef\mathcal I_{1,1}}
 +
 \underbrace{\frac12\int_{\partial^*E_\rho}|a-1|
 \left|\nu_\rho-\frac{x-\bar z_\rho}{|x-\bar z_\rho|}\right|^2\dHn}
 _{\eqdef\mathcal I_{1,2}}.
\end{align*}

The first integral can be estimated with Gauss--Green. Consider the vector
field
\[
 X(x)\eqdef\left|\frac{|x-\bar z_\rho|}{\rho}-1\right|
 \frac{x-\bar z_\rho}{|x-\bar z_\rho|}.
\]
Its divergence is locally integrable and equals
\[
 \ddiv X(x)=
 \begin{cases}
 \dfrac n\rho-\dfrac{n-1}{|x-\bar z_\rho|},
 &|x-\bar z_\rho|>\rho,\\[8pt]
 \dfrac{n-1}{|x-\bar z_\rho|}-\dfrac n\rho,
 &|x-\bar z_\rho|<\rho.
 \end{cases}
\]
On the other hand $X\big|_{\partial B_\rho}=0$, so that
\[
  \int_{\partial B_\rho}X\cdot\nu_{B_\rho}\dHn=\int_{B_\rho}\ddiv X\dd x=0,
\]
allowing us to write
\begin{equation}\label{eq.I11_estimate}
  \begin{aligned}
    \mathcal{I}_{1,1}&=
    \int_{E_\rho}\ddiv X(x)\dd x-\int_{B_\rho}\ddiv X\dd x \\
    &=\int_{E_\rho\setminus B_\rho}\ddiv X(x)\dd x-\int_{B_\rho\setminus E_\rho}\ddiv X\dd x\\
    &=\int_{E_\rho\setminus B_\rho}\left(\frac{n}{\rho}-\frac{n-1}{|x-z|}\right)\dd x
    +\int_{B_\rho\setminus E_\rho}\left(\frac{n}{\rho}-\frac{n-1}{|x-z|}\right)\dd x\\
    &\le\frac{n}{\rho}\bigl|E_\rho\,\Delta\,B_\rho(\overline{z}_\rho)\bigr|
    \le C(n,R,\rho_*)D_I(t(\rho))^{1/2}.
  \end{aligned}
\end{equation}

For the second integral, $|x-\bar z_\rho|\le2R$ on $E_\rho$, and thus
\begin{align}
 \mathcal I_{1,2}
 &\le\frac{2R+1}{2\rho_*}
 \int_{\partial^*E_\rho}
 \left|\nu_\rho-\frac{x-\bar z_\rho}{|x-\bar z_\rho|}\right|^2\dHn\notag\\
 &\le C(n,R,\rho_*)D_I(t(\rho)).
 \label{eq.I12_estimate}
\end{align}
Combining these estimates and using $D_I(t(\rho))\le\eta_I\le1$, we obtain
\begin{align*}
 \int_{\partial^*E_\rho}
 |H_{\bar z_\rho,\rho}-\overline V_\rho|\dHn
 &\le C(n,R,\rho_*)
 \bigl(D_I(t(\rho))^{1/2}+D_I(t(\rho))\bigr)\\
 &\le C(n,R,\rho_*)D_I(t(\rho))^{1/2}.
\end{align*}
Squaring proves the first integral estimate in item~(3).

It remains to compare the centers $z_\rho$ and $\bar z_\rho$. Let
$d=|z_1-z_2|$. For $0\le d\le2\rho$, the exact lens formula gives
\[
 |B_\rho(z_1)\Delta B_\rho(z_2)|
 =4\omega_{n-1}\int_0^{d/2}(\rho^2-s^2)^{(n-1)/2}\dd s.
\]
Using this formula for $d\le\rho$, and monotonicity for $d\ge\rho$, yields
\begin{equation}\label{eq.lens_lower_bound}
 |B_\rho(z_1)\Delta B_\rho(z_2)|
 \ge c_{\rm lens}(n)\rho^{n-1}\min\{|z_1-z_2|,\rho\},
 \quad
 c_{\rm lens}(n)=2\omega_{n-1}\left(\frac34\right)^{(n-1)/2}.
\end{equation}
From the minimality of $z_\rho$ and~\eqref{eq.local_FJ_volume}, we get
\begin{align*}
 c_{\rm lens}(n)\rho^{n-1}
 \min\{|z_\rho-\bar z_\rho|,\rho\}
 &\le|B_\rho(z_\rho)\Delta B_\rho(\bar z_\rho)| 
 \le|E_\rho\Delta B_\rho(z_\rho)|
 +|E_\rho\Delta B_\rho(\bar z_\rho)|\\
 &\le2|E_\rho\Delta B_\rho(\bar z_\rho)| 
 \le C(n,R,\rho_*)D_I(t(\rho))^{1/2}.
\end{align*}
Reducing the computable threshold $\eta_I(n,R,\rho_*)$ if necessary, the
minimum is attained by its first branch. Hence
\[
 |z_\rho-\bar z_\rho|
 \le C(n,R,\rho_*)D_I(t(\rho))^{1/2}.
\]
After one further reduction of $\eta_I$, the identity
$\Per(E_\rho)^2=\Per(B_\rho)^2+D_I(t(\rho))$ also gives
$\Per(E_\rho)\le2\Per(B_\rho)$. Consequently,
\begin{align*}
 \int_{\partial^*E_\rho}
 |H_{z_\rho,\rho}-H_{\bar z_\rho,\rho}|\dHn
 &\le\frac{\Per(E_\rho)}{\rho}|z_\rho-\bar z_\rho| \le C(n,R,\rho_*)D_I(t(\rho))^{1/2}.
\end{align*}
Finally, another application of the triangle inequality gives
\begin{align*}
 \int_{\partial^*E_\rho}|H_{z_\rho,\rho}-\overline V_\rho|\dHn
 &\le\int_{\partial^*E_\rho}|H_{z_\rho,\rho}-H_{\bar z_\rho,\rho}|\dHn 
 +
 \int_{\partial^*E_\rho}|H_{\bar z_\rho,\rho}-\overline V_\rho|\dHn\\
 &\le C(n,R,\rho_*)D_I(t(\rho))^{1/2}.
\end{align*}
Squaring finishes the proof of item~(3).
\end{proof}

\end{document}
